\clearpage
\fancyhead[L]{\small DIAMETER OF THE SUN.}
\hypertarget{native-524-lower}{}\label{native-524-lower}
\begin{center}\small Göttingische gelehrte Anzeigen. No. 113. pp. 1121..1126. 17 July 1809.\\[3pt]\rule{.72\linewidth}{.3pt}\end{center}

In a letter of 8 June to our Professor \textsc{Gauss}, Mr \textsc{von Lindenau} reported the \emph{results of investigations into the diameter of the Sun}, which deserve special attention. He was prompted to make these investigations when he tried to determine the Sun’s diameter from the duration of its transits, using a fairly substantial number of culminations of the Sun observed at the Seeberg Observatory. He found differences that he hesitated to ascribe solely to observational errors, since they seemed very clearly to follow a definite \emph{period}. He therefore decided to examine all the Greenwich observations from 1750 to 1786 for this purpose, and the results are indeed very remarkable. Mr \textsc{von Lindenau} first grouped the results by individual year for the diameters of the Sun, all reduced to the mean distance of the Earth from the Sun, as the following table shows:

\SourceLeaf{525}\begin{center}\small
\setlength{\tabcolsep}{3pt}
\begin{tabular}{c|c|c||c|c|c}
\shortstack{Year of\\obser-\\vation}&\shortstack{Mean result for\\the solar\\semidiameter at\\the mean\\distance}&\shortstack{Number of\\obser-\\vations}&\shortstack{Year of\\obser-\\vation}&\shortstack{Mean result for\\the solar\\semidiameter at\\the mean\\distance}&\shortstack{Number of\\obser-\\vations}\\ \hline
1750 & $962''63$ & 37 & 1773 & $961''42$ & 56 \\
1751 & $961.82$ & 42 & 1774 & $961.22$ & 32 \\
1752 & $961.61$ & 61 & 1775 & $960.86$ & 62 \\
1753 & $961.81$ & 53 & 1776 & $960.75$ & 53 \\
1754 & $962.30$ & 71 & 1777 & $960.45$ & 81 \\
1755 & $961.81$ & 46 & 1778 & $960.36$ & 71 \\
1765 & $962.44$ & 65 & 1779 & $960.92$ & 41 \\
1766 & $962.70$ & 46 & 1780 & $960.80$ & 37 \\
1767 & $961.57$ & 77 & 1781 & $960.09$ & 50 \\
1768 & $961.53$ & 72 & 1782 & $960.26$ & 37 \\
1769 & $961.23$ & 65 & 1783 & $959.84$ & 42 \\
1770 & $961.66$ & 47 & 1785 & $959.93$ & 65 \\
1771 & $961.82$ & 64 & 1786 & $959.65$ & 55 \\
1772 & $961.84$ & 56 & & & \\
\end{tabular}\end{center}

It is very striking that the Sun’s semidiameter decreases steadily here, step by step. Yet so considerable a decrease in the body of the Sun itself in so short a time is quite inconceivable. A decrease of 3 seconds in its semidiameter would amount to nearly a hundredth of its entire volume. The cause of this phenomenon can therefore presumably be sought only in the telescope, in the observer’s eye, or in the way the observations were made. — Grouping the individual determinations by season gave equally striking results. Throughout this whole series of years, with only a few exceptions, there were unmistakable changes that recurred periodically. The mean values for the individual months, all reduced to the 15th, were as follows:

\begin{center}\begin{tabular}{l|c||l|c}
&\shortstack{Semidiameter\\of the Sun}&&\shortstack{Semidiameter\\of the Sun}\\ \hline
January & $960''17$ & July & $960''14$ \\
February & $961.16$ & August & $961.00$ \\
March & $961.52$ & September & $961.70$ \\
April & $961.22$ & October & $961.80$ \\
May & $961.20$ & November & $961.16$ \\
June & $960.00$ & December & $960.43$ \\
\end{tabular}\end{center}

There are two maxima, which can be placed at the beginning of April and October, and two minima, which can be placed at the beginning of January and July. The difference between the extremes can be estimated at one and a half seconds for the semidiameter and thus three seconds for the diameter. To some it may admittedly seem rather doubtful that such small quantities can be settled by \emph{timing observations} that, taken individually, often show much larger differences. On the other hand, the great accuracy with which \textsc{Bradley} and \textsc{Maskelyne} always observed at Greenwich, the extraordinarily large number of observations, and the regularity of the changes in the results all weigh too heavily against treating these differences as mere chance and as consequences of observational errors. One might first think of attributing them to the influence of the seasons on the atmosphere, and thus to a change in irradiation, or to an effect on the sensitivity of the eye. But in that case the two different extremes would, to all appearances, have to fall in the hottest and coldest periods, whereas the observa\SourceLeaf{518}tions give the same diameter at both times. An \emph{elliptical shape of the solar body}, on the other hand, would account for the phenomenon very naturally. Suppose the Sun is an ellipsoid generated by revolution about its axis of rotation. Denote the radius of the solar equator by $a$ and the semiaxis by $b$. Further, denote the inclination of the solar equator to the Earth’s equator by $i$, the right ascension of the ascending node of the former plane on the latter by $N$, and the right ascension of the Sun by $A$. Finally, denote the apparent semidiameter of a sphere of radius $a$ seen from distance $1$ by $r$. A simple derivation shows that the semidiameter deduced from the observed culmination time and reduced to the mean distance will be

\[ r\left(1-\tfrac{1}{2}\left(1-\frac{bb}{aa}\right)\cos(A-N)^2\sin i^2\right) \]

Earlier astronomers determined the position of the solar equator. Taking its ascending node on the ecliptic to be at longitude $78^\circ$ and its inclination to the ecliptic to be $7\tfrac{1}{2}$ degrees gives $i = 26^\circ \, 2'$, $N = 16^\circ \, 55'$. The extreme values of the semidiameter will therefore occur when $A-N$ is either $0$, or $90^\circ$, or $180^\circ$, or $270^\circ$. In the first and third cases, on 8 April and 12 October, that semidiameter will be $= r\left(1-\left(1-\dfrac{bb}{aa}\right).\,0{,}09631\right)$. In the second and fourth cases, on 8 July and 6 January, the semidiameter will be found to be $= r$. The former will be smaller or larger according as the Sun is a flattened or elongated spheroid. It is very curious that the observed phenomena require the latter hypothesis under this explanation. In fact, they are represented very well if one sets $\dfrac{bb}{aa}-1 = \tfrac{1}{58}$, or $\dfrac{b}{a} = \tfrac{17}{116}$, or sets the equatorial flattening of the solar body at $\tfrac{1}{117}$. Mr \textsc{von Lindenau} finds somewhat less, namely $\tfrac{1}{140}$. — To see what direct measurements would establish about this shape, Mr \textsc{von Lindenau} also derived the vertical solar diameters from the differences between the zenith distances of the two limbs, using just as many observations. The results are as follows:

\begin{center}\begin{tabular}{c|c||c|c}
Year & \shortstack{Mean solar\\semidiameter} & Year & \shortstack{Mean solar\\semidiameter} \\ \hline
1765 & $963''81$ & 1775 & $963''15$ \\
1766 & $964.28$ & 1776 & $962.71$ \\
1767 & $963.88$ & 1777 & $961.74$ \\
1768 & $963.40$ & 1778 & $962.47$ \\
1769 & $964.07$ & 1779 & $962.85$ \\
1770 & $963.31$ & 1780 & $962.63$ \\
1771 & $962.89$ & 1781 & $961.94$ \\
1772 & $963.36$ & 1782 & $961.84$ \\
1773 & $963.16$ & 1783 & $961.52$ \\
1774 & $963.46$ & & \\
\end{tabular}\end{center}

Here too, as with the horizontal diameters, there is a progressive decrease. The vertical diameters are consistently greater than the horizontal ones by a few seconds ($3''65$ on average). Mr \textsc{von Lindenau} calculated an equatorial flattening of $\tfrac{1}{288}$ from this. Professor \textsc{Gauss} finds somewhat less still, namely $\tfrac{1}{420}$. Thus both the comparison of the horizontal solar diameters measured at different seasons with one another and the comparison of the mean horizontal diameters with the mean vertical ones point to an elongated shape for the Sun, insofar as the precision of the measurements themselves can be trusted, even though the estimates of ellipticity do not agree. The harder it is to understand how the surface of the solar body could be in equilibrium with such a shape, the more eager \SourceLeaf{519}one will be to learn the details of Mr \textsc{von Lindenau}’s interesting investigation soon.

We also take this opportunity to share a list that Mr \textsc{von Wisniewski}, an extraordinary member of the St Petersburg Academy of Sciences, sent from Astracan to Dr \textsc{Olbers}, who passed it on to us. The list gives several occultations of stars by the Moon that are still to occur this year and that are absent from \textsc{Bode}’s Astronomisches Jahrbuch. We wish to draw observers’ attention to them and thereby perhaps obtain some observations corresponding to those Mr \textsc{von Wisniewski} is to make in the Russian Empire. The list is calculated for the horizon of a place at latitude $52^{\circ}\,31'$ and longitude $2^{\mathrm u}\,0^{\mathrm m}$ east of Greenwich. — —

Several of these occultations are already listed in the Connoissance des tems.
